%============================================================
% MTH 112 Project - Sum and Difference Identities
% Updated 202504
%============================================================

\section{Sum and Difference Identities} \label{section-sum-and-difference-identities}
In this section, we will learn to use sum and difference formulas for cosine, tangent, and to use sum and difference formulas to verify identities.

\textbook{\href{https://openstax.org/books/algebra-and-trigonometry-2e/pages/9-2-sum-and-difference-identities}{9.2}}

\subsection*{Preparation Exercises} \label{preparation-exercises-section-sum-and-difference-identities}

\begin{myPrep}
Answer the following without using a calculator. These questions are intended to check the prerequisite skills needed to complete the rest of the lab.
	\begin{enumerate}
		\item Write down the memorized values from the unit circle.
			\begin{enumerate}[itemsep=5mm]
				\begin{multicols}{3}
					\item $\sin\left(\pi\right)$
					\vfill
					\item $\sin\left(\frac{\pi}{3}\right)$
					\vfill
					\item $\sin\left(\frac{5\pi}{3}\right)$
					\vfill
					\item $\sin\left(\frac{3\pi}{2}\right)$
					\vfill
					\item $\sin\left(\frac{7\pi}{6}\right)$
					\vfill
					\item $\sin\left(-\frac{3\pi}{4}\right)$
					\vfill
				\end{multicols}
			\end{enumerate}
			\vfill
		\item Count from $0$ to $\pi$ by multiples of $\frac{\pi}{12}$, and then reduce those fractions. Here's how to start out: $\frac{0\pi}{12}$, $\frac{\pi}{12}$, $\frac{2\pi}{12}$, $\frac{3\pi}{12}$, \ldots
		\vfill
		\item Write down the three Pythagorean Identities.
		\vfill
		\item True or False: $\sin(45^\circ+60^\circ)=\sin(45^\circ)+\sin(60^\circ)$. Explain your reasoning.
		\vfill
		\item Verify the identity $\frac{\sec(\phi)\sin(\phi)}{\tan(\phi)+\cot(\phi)}=\sin^2(\phi)$.
		\vfill
	\end{enumerate}
\end{myPrep}
\vfill

%======================================================
\newpage
%======================================================

\subsection*{Practice Exercises} \label{practice-exercises-section-sum-and-difference-identities}

\begin{myPractice}
	Use the sum and difference formulas and your memorized unit circle values to evaluate the following expressions without a calculator.
	\begin{enumerate}
		\begin{multicols}{2}
			\item $\sin(60^\circ-45^\circ)$
			\vspace{3cm}
			\item $\sin\left(\frac{2\pi}{3}+\frac{\pi}{4}\right)$
			\vspace{3cm}
			\item $\cos(105^\circ)$
			\item $\cos\left(\frac{17\pi}{12}\right)$
			\vspace{3cm}
			\item $\tan(75^\circ)$
			\vspace{3cm}
			\item $\tan\left(\frac{\pi}{3}-\frac{\pi}{4}\right)$
		\end{multicols}
	\end{enumerate}
\end{myPractice}
\vspace{3cm}

\begin{myPractice}
	Imagine that there are two angles, $\alpha$ and $\beta$, not in the same triangle, such that $\sin(\alpha)=\frac{2}{7}$ where $\frac{\pi}{2}<\alpha<\pi$ and $\cos(\beta)=-\frac{5}{6}$ where $\pi<\beta<\frac{3\pi}{2}$.
	\begin{enumerate}
		\item Find the value of $\cos(\alpha)$.
		\vfill

%======================================================
\newpage
%======================================================

		\item Find the value of $\sin(\beta)$.
		\vfill
		\vfill
		\item Find the value of $\sin(\alpha+\beta)$.
		\vfill
		\item Find the value of $\cos(\alpha-\beta)$.
		\vfill
	\end{enumerate}
\end{myPractice}

\begin{myPractice}
	Verify the identities algebraically.
	\begin{enumerate}
		\item $\cos(\lambda+\nu)\cos(\lambda-\nu)=\cos^2(\lambda)-\sin^2(\nu)$
		\vfill
		\vfill
		\item $\tan\left(\rho+\frac{\pi}{4}\right)=\frac{\cos(\rho)+\sin(\rho)}{\cos(\rho)-\sin(\rho)}$
		\vfill
		\vfill
	\end{enumerate}
\end{myPractice}

%======================================================
\newpage
%======================================================

\subsection*{Definitions} \label{definitions-sum-and-difference-identities}

\begin{myDefinition}[{Sum and Difference Identities}]~\\[0.5mm]
There are \textbf{sum and difference identities} for each of the trigonometric functions, but we will only focus on those for sine, cosine, and tangent.
	\begin{itemize}
		\item $\cos(\alpha+\beta)=\cos(\alpha)\cos(\beta)-\sin(\alpha)\sin(\beta)$
		\vfill
		\item $\cos(\alpha-\beta)=\cos(\alpha)\cos(\beta)+\sin(\alpha)\sin(\beta)$
		\vfill
		\item $\sin(\alpha+\beta)=\sin(\alpha)\cos(\beta)+\cos(\alpha)\sin(\beta)$
		\vfill
		\item $\sin(\alpha-\beta)=\sin(\alpha)\cos(\beta)-\cos(\alpha)\sin(\beta)$
		\vfill
		\item $\tan(\alpha+\beta)=\frac{\tan(\alpha)+\tan(\beta)}{1-\tan(\alpha)\tan(\beta)}$
		\vfill
		\item $\tan(\alpha-\beta)=\frac{\tan(\alpha)-\tan(\beta)}{1+\tan(\alpha)\tan(\beta)}$
		\vfill
	\end{itemize}
\end{myDefinition}


%======================================================
\newpage
%======================================================

\subsection*{Exit Exercises} \label{exit-exercises-section-sum-and-difference-identities}

\begin{myExit}
	Evaluate the expression $\cos(\ang{285})$ without a calculator.
\end{myExit}
\vfill
% \begin{myExit}
% 	Evaluate the expression $\sin\left(\frac{7\pi}{12}\right)$ without a calculator.
% \end{myExit}

\begin{myExit}
	Imagine that there are two angles, $\kappa$ and $\tau$, not necessarily in the same triangle, such that $\sin(\kappa)=-\frac{3}{4}$ where $\frac{3\pi}{2}<\kappa<2\pi$ and $\cos(\tau)=\frac{3}{8}$ where $0<\tau<\frac{\pi}{2}$. 
		\begin{enumerate}
			\item Find the values of $\cos(\kappa)$ and $\sin(\tau)$.
			\vfill
			\item Find the value of $\sin(\kappa-\tau)$.
			\vfill
			\item Find the value of $\tan(\kappa+\tau)$.
			\vfill
		\end{enumerate}
\end{myExit}
\exitlikert{sum and difference identities}